\section{Ordered and reversed charges: the limit on the cohomology}\label{sec:convergence}

By Theorem~\ref{thm:cancellation} the long-time behaviour is that of the
metrics induced on the cohomology representation, which move by its own
phases up to an exponentially small error. We treat ordered charges with
semistable and with polystable cohomology, and reversed charges.

For ordered charges, the \emph{Kempf--Ness function} of the cohomology
representation for the stability parameter $(\theta_0,\theta_1)$ of
\eqref{eq:theta} is the function of a pair $(h_0,h_1)$ of Hermitian metrics
on $\ker D$ and $W$
\begin{equation}\label{eq:KN}
 \cK_{\coh}(h_0,h_1)=q+\theta_0\log\det h_0-\theta_1\log\det h_1,\qquad
 q=\sum_j\tr\bigl(h_0^{-1}b_j^*h_1b_j\bigr),
\end{equation}
where $b_j$ is the matrix of $A_j|_{\ker D}$ in fixed reference bases,
$b_j^*$ its conjugate transpose, and $h_0,h_1$ stand for the Gram
matrices. Writing $h_i=g_i^*g_i$, one has
$\cK_{\coh}=\norm{g\cdot b}^2+\theta_0\log|\det g_0|^2-\theta_1\log|\det g_1|^2$
with $b=(b_j)$ and $(g\cdot b)_j=g_1b_jg_0^{-1}$; so $\cK_{\coh}$ is the
Kempf--Ness function of the point $b$ on the symmetric space
$\mathrm{GL}_{n_0}\times\mathrm{GL}_{n_1}/\mathrm U(n_0)\times\mathrm U(n_1)$
of pairs of metrics. It is invariant under a common rescaling of $h_0,h_1$ because
$n_0\theta_0=n_1\theta_1$, and its critical points are the zeros of the
moment map $\mu=(C_0-\theta_0\id,\,C_1-\theta_1\id)$
(Lemma~\ref{lem:knvariation}).

Put $\cK_*:=\inf\cK_{\coh}\in[-\infty,\infty)$. A zero of the moment map exists
if and only if the cohomology representation is polystable
\cite[Proposition~6.5]{kin94}, and $\cK_*>-\infty$ if and only if it is
semistable \cite[Theorem~5.1]{hkkp23}. King's results are stated for integral
characters. They apply here because $(\theta_0,\theta_1)$ is a positive
multiple $c\,(n_1,n_0)$ of an integral vector, and the Kempf--Ness function
for the stability parameter $c\,(n_1,n_0)$ at the arrows $b$ is $c$ times
the one for $(n_1,n_0)$ at the arrows $b/\sqrt c$, which have the same
subrepresentations.

For the second statement we recall the argument. If a subrepresentation
$(V',W')$ of dimension vector $(a_0,a_1)$ has $a_0\theta_0-a_1\theta_1>0$, let
$h_0=1\oplus e^\tau$ on $V'\oplus V'^\perp$ and $h_1=1\oplus e^\tau$ on
$W'\oplus W'^\perp$, for reference metrics; then $q$ stays bounded as
$\tau\to\infty$, and
$\cK_{\coh}=q-\tau(a_0\theta_0-a_1\theta_1)\to-\infty$, since
$n_0\theta_0=n_1\theta_1$. Conversely, for a semistable representation and
an integral character, the closure of the orbit of a lift of the
representation to the total space of the character does not meet the zero
section \cite[Proposition~3.1 and Lemma~2.2]{kin94}, the norm function on
the total space considered in \cite[\S6]{kin94} is proper on this closed set
\cite[Lemma~6.3]{kin94}, hence bounded below by a
positive constant there, and $\frac12\cK_{\coh}$ is its logarithm on the
orbit.

We write
$\norm\mu^2=\norm{\mu_0}_{\HS}^2+\norm{\mu_1}_{\HS}^2$ and
$\norm\delta^2=\norm{\delta_0}_{\HS}^2+\norm{\delta_1}_{\HS}^2$; for a pair of
matrices $\norm\cdot$ denotes this Hilbert--Schmidt norm, and for a single
matrix it remains the operator norm.

On a representation, regarded as a complex in degree zero, the normalized
metrics move with velocity $(G_0(C_0),-G_1(C_1))$ for increasing functions
$G_i$ with $G_i(\theta_i)=0$ (see \eqref{eq:velocities}), whereas the downward
gradient flow of the Kempf--Ness function has velocity
$(C_0-\theta_0,-(C_1-\theta_1))$. So the phase flow is not that gradient flow,
but it has the same stationary points and its velocity pairs negatively with
the gradient $(-\mu_0,\mu_1)$.

\subsection{The moment map tends to zero}
Let the charges be ordered. Besides \eqref{eq:constants} put
\[
 \Gamma=\max\Bigl\{\frac{(|z_0|+\bar N)^2}{y_0},\frac{(|z_1|+\bar N)^2}{y_1}\Bigr\},\qquad
 c_\delta^2=\frac{2\pi\alpha_1s_1}{y_0}+\frac{s_1^2}{y_1^2}
\]
(the formula for $c_\delta$ makes sense for all charges and is used again in
Theorem~\ref{thm:reversed}). We note that
\begin{equation}\label{eq:thetabound}
 \theta_0,\ \theta_1\ \leq\ \vartheta\ \leq\ \bar N .
\end{equation}
Indeed $\theta_1=x_1-y_1\cot\varphi$. If $\cot\varphi\geq0$, then
$0<\theta_1\leq x_1\leq M_1/\pi$. If $\cot\varphi<0$, then
$\frac\pi2<\varphi<\alpha_0$, so
$|\cot\varphi|=\cot(\pi-\varphi)\leq(\pi-\varphi)^{-1}<(\pi-\alpha_0)^{-1}\leq\rho^{-1}$
and $\theta_1\leq\max(x_1,0)+y_1/\rho\leq M_1/\pi+M_1/(2\rho)$. In both cases
$\theta_1\leq2M_1/\rho$, that is, $\vartheta=n_1\theta_1\leq\bar N$. Let
\[
 \mathfrak h(t)=\bigl(e^{\varphi t}h_V(t)|_{\ker D},\ e^{\varphi t}h_W(t)\bigr)
\]
be the pair of induced metrics on $\ker D$ and $W$, multiplied by
$e^{\varphi t}$. The function $\cK_{\coh}$ is invariant under common rescaling, so
$\mathcal E(t):=\cK_{\coh}(\mathfrak h(t))-\cK_*$ is the excess of the induced pair.

\begin{theorem}[Decay of the moment map]\label{thm:relaxation}
Let the charges be ordered and the cohomology representation semistable.
Then along every trajectory, for $t\geq t_2$,
\begin{equation}\label{eq:KNderivative}
 \frac{d}{dt}\cK_{\coh}(\mathfrak h(t))
 \leq-\frac1\Gamma\norm\mu^2+\norm\mu\norm\delta
 \leq-\frac1{2\Gamma}\norm\mu^2+\frac\Gamma2\norm\delta^2,
 \qquad\norm{\delta(t)}\leq c_\delta\,e^{-\frac{\alpha_1}4(t-t_1)},
\end{equation}
and
\[
 \mathcal E(t)+\frac1{2\Gamma}\int_{t_2}^t\norm\mu^2
 \leq\mathcal E^\sharp:=\mathcal E(0)+2\pi(N_0+\vartheta)\,t_2+\frac\Gamma{\alpha_1}c_\delta^2 ,
\]
with $N_0=\max\{N(0),n_1M_1/\rho\}$ as in Theorem~\ref{thm:energy} and
$\delta=(\delta_0,\delta_1)$, $\delta_0=\Phi_{HH}-\phi_0$,
$\delta_1=\Phi_W-\phi_1$, as in Section~\ref{sec:setting}. Moreover
$\mu(t)\to0$, $\cK_{\coh}(\mathfrak h(t))\to\cK_*$ and $m(h(t))\to|Z|$ as $t\to\infty$.
\end{theorem}

\begin{proof}
By \eqref{eq:flow} and \eqref{eq:adapted} the pair $\mathfrak h(t)$ has velocities
$\xi_0=\varphi-\Phi_{HH}$ and $\xi_1=\varphi-\Phi_W$ in the frames induced by
the adapted frames; that is,
\begin{equation}\label{eq:velocities}
 \begin{gathered}
 \xi_0=G_0(C_0)-\delta_0,\qquad\xi_1=-G_1(C_1)-\delta_1,\\
 G_0(x)=\varphi-\arg(z_0+x),\qquad G_1(x)=\arg(z_1-x)-\varphi .
 \end{gathered}
\end{equation}
The function $G_0$ is increasing with $G_0(\theta_0)=0$, because
$\Ima(e^{-\ii\varphi}(z_0+x))=c_0-\sigma x$, and
$G_0'(x)=y_0/|z_0+x|^2\in[y_0/(|z_0|+\bar N)^2,\ 1/y_0]$ on $[0,\bar N]$;
likewise $G_1$ is increasing with $G_1(\theta_1)=0$ and
$G_1'(x)=y_1/|z_1-x|^2\in[y_1/(|z_1|+\bar N)^2,\ 1/y_1]$. By
Theorem~\ref{thm:energy} the spectra of $C_0,C_1$ lie in $[0,\bar N]$ for
$t\geq t_0$, and $\theta_0,\theta_1\leq\bar N$ by \eqref{eq:thetabound}. The first variation of \eqref{eq:KN} is
$d\cK_{\coh}(\xi)=-\tr(\xi_0\mu_0)+\tr(\xi_1\mu_1)$
(Lemma~\ref{lem:knvariation}), so
\[
 \frac d{dt}\cK_{\coh}(\mathfrak h(t))
 =-\tr\bigl(\mu_0G_0(C_0)\bigr)-\tr\bigl(\mu_1G_1(C_1)\bigr)
 +\tr(\mu_0\delta_0)-\tr(\mu_1\delta_1),
\]
and $\tr(\mu_0G_0(C_0))=\sum(c-\theta_0)G_0(c)\geq\Gamma^{-1}\norm{\mu_0}_{\HS}^2$
over the eigenvalues $c$ of $C_0$, likewise for $\mu_1$. This is the first
inequality of \eqref{eq:KNderivative}, for all $t\geq t_0$; the bound on
$\norm\delta$ for $t\geq t_2$ is Theorem~\ref{thm:cancellation}(i),(iv).

Since $\norm{\xi_0},\norm{\xi_1}\leq\pi$ and
$\norm{\mu_0}_1\leq q+n_0\theta_0$, $\norm{\mu_1}_1\leq q+n_1\theta_1$, we have
$|\frac d{dt}\cK_{\coh}(\mathfrak h(t))|\leq2\pi(N_0+\vartheta)$ for all $t$, hence
$\mathcal E(t_2)\leq\mathcal E(0)+2\pi(N_0+\vartheta)t_2$. Integrating
\eqref{eq:KNderivative} from $t_2$ and using
$\frac\Gamma2\int_{t_2}^\infty\norm\delta^2\leq\frac\Gamma{\alpha_1}c_\delta^2$
and $\mathcal E\geq0$ gives the displayed bound. By \eqref{eq:adapted} and
$\norm{\Phi_\bullet}\leq\pi$,
$\norm{\dot\mu_0}_{\HS},\norm{\dot\mu_1}_{\HS}\leq2\pi q\leq2\pi N_0$, so
$\norm\mu^2$ is uniformly continuous and integrable and therefore tends to
zero. Then $C_0\to\theta_0\id$ and $C_1\to\theta_1\id$, so
$m_H\to n_0|z_0+\theta_0|+n_1|z_1-\theta_1|$. Both $z_0+\theta_0$ and
$z_1-\theta_1$ have argument $\varphi$, so this limit is
$\Rea\bigl(e^{-\ii\varphi}(Z+n_0\theta_0-n_1\theta_1)\bigr)=|Z|$, and
$m\to|Z|$ by Theorem~\ref{thm:cancellation}(iii).

It remains to show $\mathcal E(t)\to0$. By \eqref{eq:KNderivative} the
function $\mathcal E(t)+\frac\Gamma2\int_t^\infty\norm\delta^2$ is
nonincreasing on $[t_2,\infty)$, so $\mathcal E(t)$ has a limit
$\mathcal E_\infty\geq0$. Suppose $\mathcal E_\infty>0$ and choose a pair
$\mathfrak h'$ with $\cK_{\coh}(\mathfrak h')-\cK_*<\mathcal E_\infty/2$. By convexity of
$\cK_{\coh}$ along the geodesic from $\mathfrak h(t)$ to $\mathfrak h'$ in the
space of pairs of metrics, with its invariant distance $d$ (these are
recalled in the next subsection; Lemma~\ref{lem:knvariation}),
$\cK_{\coh}(\mathfrak h(t))-\cK_{\coh}(\mathfrak h')\leq\norm{\mu(t)}\,d(\mathfrak h(t),\mathfrak h')$, and the left
side is at least $\mathcal E_\infty/4$ for large $t$. By
\eqref{eq:velocities} the speed of $\mathfrak h(t)$ is at most
$\norm\mu/\min(y_0,y_1)+\norm\delta$, which is square integrable, so
$d(\mathfrak h(t),\mathfrak h')\leq c+c'\sqrt t$ by the Cauchy--Schwarz inequality. Then
$\norm{\mu(t)}^2\geq\mathcal E_\infty^2/(16(c+c'\sqrt t)^2)$ is not
integrable, a contradiction.
\end{proof}

Theorem~\ref{thm:relaxation} does not require a zero of the moment map to
exist: for semistable cohomology that is not polystable the infimum $\cK_*$
is not attained, and $\cK_{\coh}(\mathfrak h(t))\to\cK_*$ is the analogue for the phase
flow of the corresponding fact for gradient flow lines of a Kempf--Ness
function \cite[Theorem~4.3]{grs21}.

\subsection{The Kempf--Ness function on the space of metric pairs}
The convexity statements of this subsection are classical in the theory of
Kempf--Ness functions; see \cite[Theorem~4.3]{grs21}, where they are proved
for actions on closed K\"ahler manifolds. We include the short proofs for the
linear action at hand. Let $\mathcal X$ be the space of pairs $\mathfrak h=(\mathfrak h_0,\mathfrak h_1)$
of Hermitian metrics on $\ker D$ and $W$. A tangent vector at $\mathfrak h$ is a pair
$\xi=(\xi_0,\xi_1)$ of $\mathfrak h$-Hermitian endomorphisms, $\delta\mathfrak h=\mathfrak h\xi$, and
$\norm\xi^2=\tr\xi_0^2+\tr\xi_1^2$ defines an invariant Riemannian metric, for
which $\mathcal X$ is complete, simply connected and nonpositively curved.
Its geodesics are $u\mapsto \mathfrak he^{u\xi}$, the exponential map at any point is a
diffeomorphism, and closed balls are compact. Let $d$ be the distance. In a
$\mathfrak h$-orthonormal frame let $H_j$ be the matrix of $A_j|_{\ker D}$, and $C_0$,
$C_1$, $\mu_0$, $\mu_1$, $q$ as before.

\begin{lemma}\label{lem:knvariation}
$d\cK_{\coh}(\xi)=-\tr(\xi_0\mu_0)+\tr(\xi_1\mu_1)$, so the gradient of
$\cK_{\coh}$ is $(-\mu_0,\mu_1)$, of norm $\norm\mu$. Along a geodesic,
\[
 \frac{d^2}{du^2}\cK_{\coh}(\mathfrak he^{u\xi})
 =\sum_j\norm{\xi_1H_j(u)-H_j(u)\xi_0}_{\HS}^2\geq0,\qquad
 H_j(u)=e^{u\xi_1/2}H_je^{-u\xi_0/2}.
\]
\end{lemma}

\begin{proof}
In the $\mathfrak h$-orthonormal frame,
$q(\mathfrak he^{u\xi})=\sum_j\tr(e^{-u\xi_0}H_j^\dagger e^{u\xi_1}H_j)$, and
$\log\det$ is affine along geodesics with derivative $\tr\xi$. Differentiate
once and twice.
\end{proof}

So $\cK_{\coh}$ is geodesically convex and its critical points, the zeros of
the moment map, are its minimizers. Let $\mathcal Z\subset\mathcal X$ be the
set of zeros of the moment map. For $\mathfrak h\in\mathcal X$ let $\mathfrak a(\mathfrak h)$ be
the space of $\mathfrak h$-Hermitian pairs $\xi$ with $\xi_1A_j=A_j\xi_0$ on $\ker D$
for all $j$, the Hermitian infinitesimal automorphisms of the
representation.

\begin{lemma}\label{lem:minimumset}
For $p\in\mathcal Z$, $\mathcal Z=\{pe^\xi:\xi\in\mathfrak a(p)\}$. It is a
closed totally geodesic submanifold with $T_p\mathcal Z=\mathfrak a(p)$.
\end{lemma}

\begin{proof}
If $\xi\in\mathfrak a(p)$, the second derivative of $\cK_{\coh}$ along
$u\mapsto pe^{u\xi}$ vanishes identically by Lemma~\ref{lem:knvariation}, and
the first derivative vanishes at $u=0$; so $pe^\xi\in\mathcal Z$. Conversely,
if $pe^\xi\in\mathcal Z$, the convex function
$u\mapsto\cK_{\coh}(pe^{u\xi})$ takes its minimum at $u=0$ and $u=1$, hence is
constant on $[0,1]$, and its second derivative at $0$ gives
$\xi\in\mathfrak a(p)$. $\mathcal Z=\{\mu=0\}$ is closed, and it is the image
of a linear subspace under the exponential map at each of its points.
\end{proof}

\begin{lemma}\label{lem:gap}
For $p\in\mathcal Z$ put
\[
 \nu_*:=\min\Bigl\{\sum_j\norm{\xi_1H_j-H_j\xi_0}_{\HS}^2:\
 \xi\perp\mathfrak a(p),\ \norm\xi=1\Bigr\}>0,\qquad\varrho_*:=\frac{\log2}{\sqrt2}.
\]
Then $\nu_*$ does not depend on $p\in\mathcal Z$, and for every unit
$\xi\perp\mathfrak a(p)$ and $0\leq u\leq\varrho_*$,
$\frac{d^2}{du^2}\cK_{\coh}(pe^{u\xi})\geq\nu_*/2$.
\end{lemma}

\begin{proof}
The quadratic form is positive on the orthogonal complement of its kernel
$\mathfrak a(p)$. By Lemma~\ref{lem:knvariation},
$\xi_1H_j(u)-H_j(u)\xi_0=e^{u\xi_1/2}(\xi_1H_j-H_j\xi_0)e^{-u\xi_0/2}$ has
Hilbert--Schmidt norm at least
$e^{-u(\norm{\xi_1}+\norm{\xi_0})/2}\norm{\xi_1H_j-H_j\xi_0}_{\HS}$, and
$\norm{\xi_1}+\norm{\xi_0}\leq\sqrt2$ for a unit $\xi$. So the second
derivative is at least $e^{-\sqrt2u}\nu_*$. If $p'=pe^\zeta$ with
$\zeta\in\mathfrak a(p)$, then $e^{-\zeta/2}$ is a pair of automorphisms of
the representation that maps $p$-orthonormal frames to $p'$-orthonormal
frames and preserves the matrices $H_j$; conjugation by it carries
$\mathfrak a(p)$ and its orthogonal complement isometrically to those of
$p'$.
\end{proof}

\begin{lemma}\label{lem:kngrowth}
Let the representation be polystable, so that $\mathcal Z\neq\emptyset$.
For $\mathfrak h\in\mathcal X$ put $\mathsf d=d(\mathfrak h,\mathcal Z)$ and
$\mathcal E=\cK_{\coh}(\mathfrak h)-\cK_*$. Then
\begin{enumerate}
\item[(i)] $\mathcal E\leq\norm{\mu(\mathfrak h)}\,\mathsf d$;
\item[(ii)] $\mathcal E\geq\frac{\nu_*}4\min(\mathsf d^2,\varrho_*\mathsf d)$;
\item[(iii)] if $\mathsf d\leq\varrho_*$, then $\norm{\mu(\mathfrak h)}\geq\frac{\nu_*}2\mathsf d$
 and $\mathcal E\leq\frac2{\nu_*}\norm{\mu(\mathfrak h)}^2$.
\end{enumerate}
\end{lemma}

\begin{proof}
Let $p\in\mathcal Z$ be a nearest point and $\mathfrak h=pe^{\mathsf d\xi}$ with
$\norm\xi=1$; then $\xi\perp\mathfrak a(p)$. The function
$f(u)=\cK_{\coh}(pe^{u\xi})-\cK_*$ is convex with $f(0)=f'(0)=0$.
(i) $0=f(0)\geq f(\mathsf d)-\mathsf d f'(\mathsf d)$ and $f'(\mathsf d)\leq\norm{\mu(\mathfrak h)}$.
(ii) $f''\geq\nu_*/2$ on $[0,\varrho_*]$ by Lemma~\ref{lem:gap}, so
$f(u)\geq\frac{\nu_*}4u^2$ there, and $f(u)\geq f(\varrho_*)u/\varrho_*$ for
$u\geq\varrho_*$ by convexity. (iii) $\norm{\mu(\mathfrak h)}\geq f'(\mathsf d)
=\int_0^\mathsf d f''\geq\frac{\nu_*}2\mathsf d$; combine with (i).
\end{proof}

Part (iii) is a \L ojasiewicz gradient inequality with exponent $\frac12$
for $\cK_{\coh}$ near $\mathcal Z$.


\subsection{Convergence for polystable cohomology}

\begin{theorem}[Convergence to a harmonic metric]\label{thm:convergence}
Let the charges be ordered and the cohomology representation polystable.
Then along every trajectory there is a zero $\mathfrak h_\infty$ of the moment map,
that is, a harmonic metric of the cohomology representation in the sense of
\cite{hkkp26}, with
\[
 d\bigl(\mathfrak h(t),\mathfrak h_\infty\bigr)\leq\mathfrak C\,e^{-\beta t}\qquad(t\geq0)
\]
for every rate $\beta<\frac12\min\{\nu_*/(4\Gamma),\ \alpha_1/2\}$, where the
constant $\mathfrak C$ depends only on $\beta$, the dimensions, the charges, $N(0)$, $t_2$,
$\mathcal E(0)$ and $\nu_*$. Moreover, for $t\geq t_1$,
\[
 \norm{\mu(\mathfrak h(t))}\leq2\bar N\,d(\mathfrak h(t),\mathfrak h_\infty),\qquad
 0\leq m_H(t)-|Z|\leq\frac{\norm{\mu(\mathfrak h(t))}^2}{2\min(y_0,y_1)},
\]
so that the moment map of the cohomology representation and $m(h(t))-|Z|$
tend to zero exponentially.
\end{theorem}

The convergence is as Hermitian forms on the fixed vector spaces $\ker D$ and
$W$: no time-dependent automorphism of the representation is applied. The
limit depends on the initial metric through its position on $\mathcal Z$.

\begin{proof}
By Theorem~\ref{thm:relaxation}, \eqref{eq:KNderivative} holds for
$t\geq t_2$ and $\mathcal E(t)\leq\mathcal E^\sharp$ there.

\emph{Step 1.} By Lemma~\ref{lem:kngrowth}(ii),
$d(\mathfrak h(t),\mathcal Z)\leq\mathsf d_{\max}:=\max\{\varrho_*,4\mathcal E^\sharp/(\nu_*\varrho_*)\}$
for $t\geq t_2$, and by (i), $\mathcal E(t)\leq\norm{\mu(t)}\mathsf d_{\max}$.

\emph{Step 2.} By Theorem~\ref{thm:relaxation},
$\int_{t_2}^\infty\norm\mu^2\leq2\Gamma\mathcal E^\sharp$ and
$\mathcal E(t)\leq\mathcal E(t')+\frac\Gamma2\int_{t'}^\infty\norm\delta^2$ for
$t\geq t'\geq t_2$. For $T>0$ there is $t'\in[t_2+T/2,t_2+T]$ with
$\norm{\mu(t')}^2\leq4\Gamma\mathcal E^\sharp/T$, hence for $t\geq t_2+T$
\[
 \mathcal E(t)\leq\mathsf d_{\max}\sqrt{4\Gamma\mathcal E^\sharp/T}
 +\frac\Gamma{\alpha_1}c_\delta^2e^{-\frac{\alpha_1}4T}.
\]
Choose $T$, depending only on the dimensions, the charges,
$\mathcal E^\sharp$ and $\nu_*$, so that the right side is at most
$\frac{\nu_*}4\varrho_*^2$; then $d(\mathfrak h(t),\mathcal Z)\leq\varrho_*$ for
$t\geq t_3:=t_2+T$ by Lemma~\ref{lem:kngrowth}(ii). For such $t$,
Lemma~\ref{lem:kngrowth}(iii) and \eqref{eq:KNderivative} give
\[
 \dot{\mathcal E}\leq-\frac{\nu_*}{4\Gamma}\mathcal E
 +\frac\Gamma2c_\delta^2e^{-\frac{\alpha_1}2(t-t_1)},
\]
so $\mathcal E(t)\leq\mathfrak C_2e^{-2\beta t}$ for every
$2\beta<\min\{\nu_*/(4\Gamma),\alpha_1/2\}$.

\emph{Step 3.} Integrating \eqref{eq:KNderivative} over $[t,t+1]$,
$\int_t^{t+1}\norm\mu^2\leq2\Gamma\mathcal E(t)+\Gamma^2\int_t^{t+1}\norm\delta^2
\leq\mathfrak C_3e^{-2\beta t}$. By \eqref{eq:velocities} and $|G_0'|\leq1/y_0$,
$|G_1'|\leq1/y_1$, the speed of $\mathfrak h(t)$ is
$\norm\xi\leq\norm\mu/\min(y_0,y_1)+\norm\delta$. Hence
$\int_t^\infty\norm\xi\leq\mathfrak C_4e^{-\beta t}$ for $t\geq t_3$, and by
completeness $\mathfrak h(t)$ converges to some $\mathfrak h_\infty$ with
$d(\mathfrak h(t),\mathfrak h_\infty)\leq\mathfrak C_4e^{-\beta t}$; $\mu(\mathfrak h_\infty)=0$ by continuity. On
$[0,t_3]$ the speed is at most $\pi\sqrt{n_0+n_1}$, because the spectra of
$\Phi_{HH}$ and $\Phi_W$ lie in $(0,\pi)$; this gives the dependence of $\mathfrak C$.

\emph{The moment map and the mass.} Let $t\geq t_1$ and let
$u\mapsto \mathfrak h_\infty e^{u\zeta}$, $0\leq u\leq1$, be the geodesic from
$\mathfrak h_\infty$ to $\mathfrak h(t)$. The function $q$ is convex along geodesics, so
$q\leq\max(\vartheta,q(\mathfrak h(t)))\leq\bar N$ on it, by
\eqref{eq:thetabound}; here $q(\mathfrak h_\infty)=\tr C_0=\vartheta$. In the frames of
Lemma~\ref{lem:knvariation},
$\frac d{du}C_0=-\frac12(\zeta_0C_0+C_0\zeta_0)+\sum_jH_j^\dagger\zeta_1H_j$
and similarly for $C_1$, so $\norm{\frac d{du}\mu}\leq2q\norm\zeta$, and
$\mu(\mathfrak h_\infty)=0$. For the mass, the functions $x\mapsto|z_0+x|$ and
$x\mapsto|z_1-x|$ are convex with second derivatives at most $1/y_0$ and
$1/y_1$, and their first-order terms at $\theta_0,\theta_1$ cancel in the sum
over the eigenvalues of $C_0$ and $C_1$ because $\tr\mu_0=\tr\mu_1$. Combine
with Theorem~\ref{thm:cancellation}(iii).
\end{proof}

\begin{remark}\label{rem:rate}
The rate is limited by the gap $\nu_*$ of the Hessian of $\cK_{\coh}$
transversally to $\mathcal Z$ and by the decay of $A_jD^+$. When the
cohomology representation is stable, $\mathfrak a$ consists of the scalar
pairs $(c,c)$.
\end{remark}

\subsection{Semistable cohomology that is not polystable}
In this case Theorem~\ref{thm:relaxation} gives $\mu\to0$, and there is no
limit metric.

\begin{proposition}\label{prop:divergence}
Let the charges be ordered and let the cohomology representation be
semistable and not polystable. Then along every trajectory
$d(\mathfrak h(t),\mathfrak h(0))\to\infty$, and for $t\geq t_2$
\[
 d\bigl(\mathfrak h(t),\mathfrak h(t_2)\bigr)
 \leq\frac{\sqrt{2\Gamma\mathcal E^\sharp}}{\min(y_0,y_1)}\,\sqrt{t-t_2}
 +\frac{4c_\delta}{\alpha_1}.
\]
\end{proposition}

\begin{proof}
Suppose $d(\mathfrak h(t_k),\mathfrak h(0))$ is bounded along a sequence
$t_k\to\infty$. Closed balls of $\mathcal X$ are compact, so a subsequence
converges to a pair $\mathfrak h_\infty$. The moment map is continuous on
$\mathcal X$ and $\mu(\mathfrak h(t))\to0$ by
Theorem~\ref{thm:relaxation}; so $\mu(\mathfrak h_\infty)=0$ and the
representation is polystable, a contradiction. By \eqref{eq:velocities} and
$0<G_i'\leq1/y_i$ on $[0,\infty)$, the speed of $\mathfrak h(t)$ is at most
$\norm\mu/\min(y_0,y_1)+\norm\delta$. By Theorem~\ref{thm:relaxation} and the
Cauchy--Schwarz inequality
$\int_{t_2}^t\norm\mu\leq(2\Gamma\mathcal E^\sharp)^{1/2}\sqrt{t-t_2}$, and
$\int_{t_2}^\infty\norm\delta\leq4c_\delta/\alpha_1$ because $t_2\geq t_1$.
\end{proof}

The bound $\sqrt t$ is not attained in the following example, the smallest
one: there the distance grows like $\log t$, and the asymptotics can be
determined for every initial metric. Let $\mathsf R$ be the representation of
the Kronecker quiver with two arrows on $\CC^2\to\CC^2$ with the arrows
\[
 b_1=\id,\qquad b_2=E_{12},\qquad E_{12}e_2=e_1,\quad E_{12}e_1=0 .
\]
Here $n_0=n_1=2$, so $\theta_0=\theta_1=:\theta$. A subrepresentation
$(V',W')$ has $V'=b_1V'\subset W'$, so $\mathsf R$ is semistable; the
lines $\CC e_1\to\CC e_1$ form a subrepresentation of dimension vector
$(1,1)$, so it is not stable; and it is not polystable, since a direct sum
of two representations of dimension vector $(1,1)$ has arrows that are
diagonal in suitable bases of the two spaces, and then $b_1^{-1}b_2=E_{12}$
would be diagonalizable. Both $\CC e_1\to\CC e_1$ and the quotient are
isomorphic to the stable representation $\CC\to\CC$ with arrows $(1,0)$. Put
\[
 \mathsf g_0=G_0'(\theta)=\frac{y_0}{|z_0+\theta|^2},\qquad
 \mathsf g_1=G_1'(\theta)=\frac{y_1}{|z_1-\theta|^2},\qquad
 \mathsf k=\frac{2\mathsf g_0\mathsf g_1}{\mathsf g_0+\mathsf g_1},\qquad
 \Lambda_\infty=(\mathsf g_0+\mathsf g_1)\,\theta,
\]
with $G_0,G_1$ as in \eqref{eq:velocities}. For a pair of metrics
$(h_0,h_1)$ on $\mathsf R$, with Hermitian forms conjugate-linear in the
first argument, let
\[
 p_1=h_0(e_1,e_1),\qquad\chi_0=\frac{h_0(e_1,e_2)}{p_1},\qquad
 p_2=h_0(e_2-\chi_0e_1,\,e_2-\chi_0e_1),
\]
and define $p_1',\chi_1,p_2'$ in the same way from $h_1$; these are the
Gram--Schmidt data of the two metrics for the flag $\CC e_1\subset\CC^2$.
The metrics are diagonal in the basis $e_1,e_2$ exactly when
$\chi_0=\chi_1=0$, and the automorphisms $\id+xE_{12}$, $x\in\CC$, of
$\mathsf R$ shift $\chi_0$ and $\chi_1$ by $x$. Put
\[
 a=\frac{p_1'}{p_1},\qquad b=\frac{p_2'}{p_2},\qquad c=\frac{p_1'}{p_2},\qquad
 \chi=\chi_1-\chi_0 .
\]

\begin{proposition}[A semistable example]\label{prop:semistable}
Let the charges be ordered. For every initial pair of metrics on
$\mathsf R$, regarded as a complex in degree zero, along the phase flow:
\begin{enumerate}
\item[(i)] $\chi$ vanishes identically or has constant argument, it tends to
 zero with $|\chi(t)|=O(e^{-\Lambda't})$ for every $\Lambda'<\Lambda_\infty$,
 and $\chi_0$, $\chi_1$ converge to a common limit $\chi_\infty$;
\item[(ii)] $c(t)=\dfrac1{\mathsf kt}+O\Bigl(\dfrac{\log t}{t^2}\Bigr)$,
 $a-\theta=-\dfrac{\mathsf g_1}{\mathsf g_0+\mathsf g_1}c+O(c^2)$ and
 $b-\theta=-\dfrac{\mathsf g_0}{\mathsf g_0+\mathsf g_1}c+O(c^2)$;
\item[(iii)] $\displaystyle\lim_{t\to\infty}t\,\norm{\mu(t)}
 =\frac{\sqrt{2(\mathsf g_0^2+\mathsf g_1^2)}}{2\mathsf g_0\mathsf g_1}$ and
 $\cK_{\coh}(\mathfrak h(t))-\cK_*=\dfrac1{\mathsf kt}+O\Bigl(\dfrac{\log t}{t^2}\Bigr)$;
\item[(iv)] there is a constant $A>0$ with
 \[
 e^{\varphi t}p_1\sim A\,t^{-1/2},\quad e^{\varphi t}p_1'\sim\theta A\,t^{-1/2},\quad
 e^{\varphi t}p_2\sim\theta\mathsf kA\,t^{1/2},\quad
 e^{\varphi t}p_2'\sim\theta^2\mathsf kA\,t^{1/2},
 \]
 and $d(\mathfrak h(t),\mathfrak h(0))=\log t+O(1)$.
\end{enumerate}
\end{proposition}

So for every initial metric the normalized metrics decay like $t^{-1/2}$ on
the subrepresentation $\CC e_1\to\CC e_1$ and grow like $t^{1/2}$ on the
quotient. These rates hold in every basis $e_1,\ e_2-xe_1$; the value
$x=\chi_\infty$ is the one for which the metrics become diagonal, in the
sense that $h_i(e_1,e_2-\chi_\infty e_1)/h_i(e_1,e_1)\to0$ for $i=0,1$. The
exponents do not depend on the charges; the constants do, through $\theta$,
$\mathsf g_0$ and $\mathsf g_1$.

\begin{proof}
\emph{Frames.} The vectors $e_1/\sqrt{p_1}$, $(e_2-\chi_0e_1)/\sqrt{p_2}$ form
an $h_0$-orthonormal frame, and $e_1/\sqrt{p_1'}$,
$(e_2-\chi_1e_1)/\sqrt{p_2'}$ an $h_1$-orthonormal frame. Since
$e_2-\chi_0e_1=(e_2-\chi_1e_1)+\chi e_1$ and $E_{12}(e_2-\chi_0e_1)=e_1$, the
arrows have the matrices
$H_1=\bigl(\begin{smallmatrix}\sqrt a&\chi\sqrt c\\0&\sqrt b\end{smallmatrix}\bigr)$
and $H_2=\bigl(\begin{smallmatrix}0&\sqrt c\\0&0\end{smallmatrix}\bigr)$, so
\[
 C_0=\begin{pmatrix}a&\chi\sqrt{ac}\\\bar\chi\sqrt{ac}&b+c+|\chi|^2c\end{pmatrix},\qquad
 C_1=\begin{pmatrix}a+c+|\chi|^2c&\chi\sqrt{bc}\\\bar\chi\sqrt{bc}&b\end{pmatrix}.
\]
Let $F_0(x)=\arg(z_0+x)=\varphi-G_0(x)$ and
$F_1(x)=\arg(z_1-x)=\varphi+G_1(x)$; the matrices of the phases in these
frames are $F_0(C_0)$ and $F_1(C_1)$.

\emph{Evolution of the data.} For a metric $h$ with phase $\Phi$,
$\frac d{dt}h(u,w)=-h(u,\Phi w)$. Writing $e_1$, $e_2$ in the frame, one finds
for the data $(p_1,p_2,\chi_0)$ of $h_0$ and the matrix $\phi=F_0(C_0)$
\[
 \tfrac d{dt}\log p_1=-\phi_{11},\qquad\tfrac d{dt}\log p_2=-\phi_{22},\qquad
 \dot\chi_0=-\sqrt{p_2/p_1}\;\phi_{12},
\]
the second because $\det h_0=p_1p_2$ has logarithmic derivative $-\tr\phi$;
and likewise for $h_1$ with $\phi=F_1(C_1)$. For a $2\times2$ Hermitian matrix
$C$ with eigenvalues $c^\pm$ and a function $f$, $f(C)-f[c^+,c^-]\,C$ is a
scalar matrix, where $f[\,\cdot\,,\cdot\,]$ is the divided difference; so
$f(C)_{12}=f[c^+,c^-]\,C_{12}$. Let $\varrho_0=-F_0[c_0^+,c_0^-]>0$ and
$\varrho_1=F_1[c_1^+,c_1^-]>0$ for the eigenvalues of $C_0$, $C_1$. Since
$\sqrt{p_2/p_1}\sqrt{ac}=a$ and $\sqrt{p_2'/p_1'}\sqrt{bc}=b$,
\begin{equation}\label{eq:chi}
 \dot\chi_0=\varrho_0\,a\,\chi,\qquad\dot\chi_1=-\varrho_1\,b\,\chi,\qquad
 \dot\chi=-(a\varrho_0+b\varrho_1)\,\chi .
\end{equation}

\emph{Limits.} Theorem~\ref{thm:relaxation} applies to the flow of a
representation. Indeed, let $\mathsf E$ be the orthogonal direct sum of
$\mathsf R$ and a complex $(V_1\to0)\to(U\to0)$ with invertible differential.
It is of the form \eqref{eq:complex}, its endomorphism \eqref{eq:coeff} is
the direct sum of those of the two summands, so the flow preserves the
orthogonal splitting and restricts to the flow of each summand; the
cohomology of $\mathsf E$ is $\mathsf R$ with its own metrics, and
$\delta=0$. Hence $C_0\to\theta$ and $C_1\to\theta$, that is, $a\to\theta$,
$b\to\theta$, $c(1+|\chi|^2)\to0$, and $\varrho_i\to\mathsf g_i$. By
\eqref{eq:chi}, $\chi$ is a real positive multiple of $\chi(0)$, $|\chi|$ is
nonincreasing, and $a\varrho_0+b\varrho_1\to\Lambda_\infty$; this gives the
decay of $\chi$. Then
$\dot\chi_0$ and $\dot\chi_1$ are integrable, so $\chi_0$, $\chi_1$ converge,
to the same limit. This is (i).

\emph{The reduced system.} Put
\[
 w_1=-G_0(a)-G_1(a+c),\qquad w_2=-G_0(b+c)-G_1(b),\qquad
 w_3=-G_0(b+c)-G_1(a+c),
\]
functions of $(a,b,c)$. If $\chi=0$, the matrices $C_0$, $C_1$ are diagonal
and $\frac d{dt}\log a=w_1$, $\frac d{dt}\log b=w_2$,
$\frac d{dt}\log c=w_3$. In general, for a $2\times2$ Hermitian matrix $C$
with eigenvalues $c^\pm$ and a smooth function $f$, Newton interpolation of
$f$ at the nodes $C_{11},c^+,c^-$ gives
\[
 f(C)=f(C_{11})+f[C_{11},c^+]\,(C-C_{11})+f[C_{11},c^+,c^-]\,(C-C_{11})(C-c^+),
\]
and the upper left entries of $C-C_{11}$ and $(C-C_{11})(C-c^+)$ are $0$ and
$|C_{12}|^2$. So $|f(C)_{11}-f(C_{11})|\leq\frac12\sup|f''|\,|C_{12}|^2$, and
likewise for the other diagonal entry. Here $|C_{12}|^2\leq|\chi|^2c\,(a+b)$,
and the diagonal entries of $C_0$, $C_1$ differ from those for $\chi=0$ by
$|\chi|^2c$. Hence, with
$\varepsilon(t):=|\chi(t)|^2$,
\begin{equation}\label{eq:reduced}
 \tfrac d{dt}\log a=w_1+O(\varepsilon c),\qquad
 \tfrac d{dt}\log b=w_2+O(\varepsilon c),\qquad
 \tfrac d{dt}\log c=w_3+O(\varepsilon c),
\end{equation}
and $\frac d{dt}\log(e^{\varphi t}p_1)=G_0(a)+O(\varepsilon c)$, and similarly
for $p_2$, $p_1'$, $p_2'$ with $G_0(b+c)$, $-G_1(a+c)$, $-G_1(b)$.

\emph{The fast variables.} Differentiating along \eqref{eq:reduced},
\[
 \dot w_1=-\Lambda_aw_1-G_1'(a+c)\,c\,w_3+O(\varepsilon c),\qquad
 \dot w_2=-\Lambda_bw_2-G_0'(b+c)\,c\,w_3+O(\varepsilon c),
\]
with $\Lambda_a=a(G_0'(a)+G_1'(a+c))$ and $\Lambda_b=b(G_0'(b+c)+G_1'(b))$,
both tending to $\Lambda_\infty$. Near $(a,b,c)=(\theta,\theta,0)$ the
functions $(w_1,w_2,c)$ are real-analytic coordinates, with
$w_1=-(\mathsf g_0+\mathsf g_1)(a-\theta)-\mathsf g_1c$ and
$w_2=-(\mathsf g_0+\mathsf g_1)(b-\theta)-\mathsf g_0c$ to first order.
Therefore, with $|w|=(w_1^2+w_2^2)^{1/2}$,
\begin{equation}\label{eq:slow}
 a-\theta=-\frac{w_1+\mathsf g_1c}{\mathsf g_0+\mathsf g_1}+O(|w|^2+c^2),\qquad
 w_3=-\mathsf kc+\frac{\mathsf g_1w_1+\mathsf g_0w_2}{\mathsf g_0+\mathsf g_1}
 +O(|w|^2+c^2),
\end{equation}
and similarly for $b$; the factor of $c$ in $w_3$ is
$-(\mathsf g_0+\mathsf g_1)+(\mathsf g_0^2+\mathsf g_1^2)/(\mathsf g_0+\mathsf g_1)=-\mathsf k$.

We claim $|w|=O(c^2)$. Choose $T$ such that for $t\geq T$:
$\Lambda_a,\Lambda_b\geq\Lambda_\infty/2$,
$|\frac d{dt}\log c|\leq\Lambda_\infty/16$, $|w|\leq1$, $c\leq1$,
\eqref{eq:slow} holds, and $a\varrho_0+b\varrho_1\geq\Lambda_\infty/2$, so
that $\varepsilon(s)\leq\varepsilon(T)e^{-\Lambda_\infty(s-T)}$ for $s\geq T$
by \eqref{eq:chi}. Then, for constants $K',K''$ and after enlarging $T$,
\[
 \frac d{dt}|w|\leq-\frac{\Lambda_\infty}4|w|+K'\,c^2+K''\,\varepsilon\,c
\]
in the sense of upper right derivatives. Since
$c(s)\leq c(t)e^{\frac{\Lambda_\infty}{16}(t-s)}$ and
$e^{-\frac{\Lambda_\infty}{16}(t-T)}\leq c(t)/c(T)$ for $T\leq s\leq t$,
\[
 |w(t)|\leq\Bigl(\frac{c(t)}{c(T)}\Bigr)^4|w(T)|+\frac{8K'}{\Lambda_\infty}\,c(t)^2
 +K''c(t)\,\varepsilon(T)\int_T^te^{-\frac{3\Lambda_\infty}{16}(t-s)}e^{-\Lambda_\infty(s-T)}ds ,
\]
and the integral is at most
$\frac{16}{13\Lambda_\infty}e^{-\frac{3\Lambda_\infty}{16}(t-T)}\leq\frac{16}{13\Lambda_\infty}(c(t)/c(T))^3$.
This proves the claim. By \eqref{eq:reduced} and \eqref{eq:slow},
$\frac d{dt}\log c=-\mathsf kc+O(c^2)+O(\varepsilon c)$, so
$\frac d{dt}(1/c)=\mathsf k+O(c)+O(\varepsilon)$; then $c=O(1/t)$,
$\int^tc=O(\log t)$ and $1/c=\mathsf kt+O(\log t)$. With \eqref{eq:slow} this
is (ii).

(iii) By (ii), up to terms $O(c^2)$ and terms of order
$|\chi|\sqrt c$, which decay exponentially,
$C_0-\theta=\frac{\mathsf g_1c}{\mathsf g_0+\mathsf g_1}\diag(-1,1)$ and
$C_1-\theta=\frac{\mathsf g_0c}{\mathsf g_0+\mathsf g_1}\diag(1,-1)$. Moreover
$\cK_{\coh}=\tr C_0+\theta\log(\det h_0/\det h_1)=a+b+c+\varepsilon c-\theta\log(ab)$,
and $x-\theta\log x$ is stationary at $\theta$, so
$\cK_{\coh}=2\theta-2\theta\log\theta+c+O(c^2)+O(\varepsilon c)$; the constant
is $\cK_*$ by Theorem~\ref{thm:relaxation}.

(iv) By \eqref{eq:reduced} and (ii),
$\frac d{dt}\log(e^{\varphi t}p_1)=\mathsf g_0(a-\theta)+O(c^2)+O(\varepsilon c)
=-\frac{\mathsf k}2c+O(c^2)+O(\varepsilon c)=-\frac1{2t}+O(\frac{\log t}{t^2})$;
the error is integrable, so $e^{\varphi t}p_1\,t^{1/2}$ has a positive limit
$A$. Then $p_1'=ap_1$, $p_2=p_1'/c$ and $p_2'=bp_2$ give the other three. The
metrics are $g_i^*\diag g_i$ with the diagonal matrices of the data and
$g_i=\bigl(\begin{smallmatrix}1&\chi_i\\0&1\end{smallmatrix}\bigr)$, which
converge. For every pair $(g_0,g_1)$ of invertible matrices the map
$(h_0,h_1)\mapsto(g_0^*h_0g_0,\,g_1^*h_1g_1)$ is an isometry of the space of
pairs of metrics, and the $g_i(t)$ stay in a compact set; so the distance
from $\mathfrak h(0)$ is
$\bigl(\sum_i\log^2(e^{\varphi t}p_i)+\sum_i\log^2(e^{\varphi t}p_i')\bigr)^{1/2}+O(1)
=\log t+O(1)$.
\end{proof}

\begin{remark}\label{rem:semistable}
Apart from the limits $a,b\to\theta$ and $c\to0$, the proof uses of
$G_0,G_1$ only that they are smooth and increasing, with positive
derivatives at $\theta$, and vanish at $\theta$. For the downward gradient
flow of the Kempf--Ness function, for which these limits hold as well, the
same computation applies with $G_i(x)=x-\theta$, and gives the same exponents
with $\mathsf g_0=\mathsf g_1=\mathsf k=1$; asymptotics of such gradient flows
at semistable points are the subject of \cite{hkkp23}. For $\chi(0)\neq0$ the rate in (i) is
not attained for $\Lambda'=\Lambda_\infty$: the divided differences satisfy
$\varrho_i=\mathsf g_i+O(c^2)$, hence
$a\varrho_0+b\varrho_1=\Lambda_\infty-\mathsf kc+O(c^2)$ by (ii), and $|\chi(t)|$
is asymptotic to a constant multiple of $t\,e^{-\Lambda_\infty t}$. For
general semistable cohomology the induced metrics are determined up to a
bounded factor by \cite[Theorem~6.3 and Corollary~6.5]{wue26b}; whether the
factor converges is open.
\end{remark}

\subsection{Reversed charges}

\begin{theorem}[Reversed charges]\label{thm:reversed}
Let $\alpha_0<\alpha_1$. Along every trajectory, for
$t\geq t_{\mathrm{rev}}:=\max\bigl\{t_2,\ t_1+\frac4{\alpha_1}\log^+\frac{2\sqrt2\,c_\delta}{\alpha_1-\alpha_0}\bigr\}$,
\[
 q(t)\leq\bar N\,e^{-\frac12(\alpha_1-\alpha_0)(t-t_{\mathrm{rev}})} ,
\]
and $m(h(t))\to n_0|z_0|+n_1|z_1|$ as $t\to\infty$. The phases of the induced
metrics on $\ker D$ and $W$ converge to $\alpha_0$ and $\alpha_1$.
\end{theorem}

\begin{proof}
By \eqref{eq:adapted},
$\dot q=\tr(C_0\Phi_{HH})-\tr(C_1\Phi_W)
=\tr(C_0\phi_0)-\tr(C_1\phi_1)+\tr(C_0\delta_0)-\tr(C_1\delta_1)$. Since
$C_0,C_1\geq0$, $\phi_0=\arg(z_0\id+C_0)\leq\alpha_0$ and
$\phi_1=\arg(z_1\id-C_1)\geq\alpha_1$. Hence, for all charges,
$\dot q\leq-(\alpha_1-\alpha_0-\norm{\delta_0}-\norm{\delta_1})\,q$. By
Theorem~\ref{thm:cancellation}(i),(iv),
$\norm{\delta_0}+\norm{\delta_1}\leq\sqrt2\,c_\delta e^{-\frac{\alpha_1}4(t-t_1)}$
for $t\geq t_2$, which is at most $\frac12(\alpha_1-\alpha_0)$ for
$t\geq t_{\mathrm{rev}}$; and $q(t_{\mathrm{rev}})\leq\bar N$. Then $C_0,C_1\to0$, so
$\phi_0\to\alpha_0$, $\phi_1\to\alpha_1$, $\delta\to0$,
$m_H\to n_0|z_0|+n_1|z_1|$, and $|m-m_H|\to0$ by
Theorem~\ref{thm:cancellation}(iii).
\end{proof}

For reversed charges every representation with nonzero arrows is unstable;
its Harder--Narasimhan factors are the two semisimple representations
supported at one vertex each, of phases $\alpha_1>\alpha_0$, and Theorem~\ref{thm:reversed} says that the flow
converges to their direct sum. By the triangle inequality
$n_0|z_0|+n_1|z_1|>|Z|$ unless $z_0,z_1$ are parallel.

\begin{remark}\label{rem:unstable}
For ordered charges and cohomology that is not semistable,
Theorem~\ref{thm:cancellation} still gives the cancellation of the acyclic
part, with arrows bounded by $q\leq\bar N$ after the time $t_0$, and reduces
the long-time behaviour
to that of the phase flow of the cohomology representation, perturbed by the
exponentially decaying $\delta_0,\delta_1$. In \cite[\S4.1]{hkkp26} the
asymptotics of the flow on an object without harmonic metric are expected to
be controlled by a refinement of its Harder--Narasimhan filtration. For the
complexes of this paper this is proved in the following form
\cite[Theorem~6.11 and Corollary~6.5]{wue26b}: the induced metrics decay at
the exponential rates given by the phases of the Harder--Narasimhan
filtration of the cohomology representation, on the steps of that
filtration, and the metric induced on each factor is determined, up to a
bounded factor, by the theory for semistable representations. The simplest case is that in which all arrows of the
cohomology representation vanish: then $H_j=0$ for all time,
$\Omega_V$ is block diagonal with the block $z_0\id$ on $\ker D$, so
$\Phi_{HH}=\alpha_0$ exactly, and $\Phi_W=\arg(z_1\id-Y)\to\alpha_1$.
\end{remark}
